% 基于深度强化学习的以太坊交易序列决策排序方法研究
% 期刊论文 LaTeX 源文件
% 编译方式：xelatex main.tex && bibtex main && xelatex main.tex && xelatex main.tex

\documentclass[a4paper, 11pt]{ctexart}

% ========== 宏包 ==========
\usepackage[left=2.5cm, right=2.5cm, top=2.8cm, bottom=2.5cm]{geometry}
\usepackage{amsmath, amssymb, amsfonts}
\usepackage{graphicx}
\usepackage{booktabs}
\usepackage{multirow}
\usepackage{algorithm}
\usepackage{algorithmic}
\usepackage{hyperref}
\usepackage{xcolor}
\usepackage{float}
\usepackage{enumitem}
\usepackage{cite}
\usepackage{url}

% ========== 算法环境中文适配 ==========
\renewcommand{\algorithmicrequire}{\textbf{输入:}}
\renewcommand{\algorithmicensure}{\textbf{输出:}}

% ========== 占位标记命令 ==========
\newcommand{\TODO}[1]{\textcolor{red}{[TODO: #1]}}

% ========== 文档信息 ==========
\title{基于深度强化学习的以太坊交易序列决策排序方法研究}
\author{
  陈庆贺 \quad 张洪豪 \\
  {\small 天津理工大学 计算机科学与工程学院，天津 300384}
}
\date{}

\begin{document}

\maketitle

% 摘要
% chapters/journal/00-abstract.tex
% 摘要

\begin{abstract}
区块链系统中，交易在区块内的执行顺序直接影响区块构建者的手续费收益、用户的确认等待体验以及系统面临的最大可提取价值（MEV）风险。现有排序方法以 Gas 手续费贪心排列为主，在追求短期收益的同时忽略了等待时间均衡性与位置敏感的安全风险；已有强化学习方法多将动作空间定义为在少量预设规则之间切换，无法实现逐笔交易粒度的排序决策。本文将区块构建过程建模为逐笔交易选择的马尔可夫决策过程，采用近端策略优化算法训练排序策略网络，奖励函数采用单步奖励、终止奖励和有效性约束的分层设计，统一建模手续费收益、等待服务、终止公平、位置敏感风险、Gas 利用和非法/提前终止行为，并通过合法性掩码机制保证动作可行性。实验基于统一主场景（$N{=}300$，$p_{\text{risk}}{=}15\%$，$\kappa{=}2.0$）的可配置仿真环境，与七类基线在共享评估池上进行对比，其中含参数基线在固定验证池上调参，正式统计以独立训练 seed 为单位。现阶段结果支持本文方法在收益与综合指标上的折中优势，但公平性与风险指标并非同时最优，约束评估也显示可行率与可行子集收益之间存在权衡。本文的实验结论限定在以太坊生态下的抽象区块构建/L2 Sequencer 排序场景，不直接等价于已完成 L1 PBS/Builder 机制级验证。

\noindent\textbf{关键词：} 区块链；交易排序；深度强化学习；序列决策；区块构建
\end{abstract}


% 第1章 引言
% chapters/journal/01-introduction.tex
% 第1章 引言

\section{引言}

区块链系统中，区块构建者负责从交易池中选取待处理交易并确定其在区块内的执行顺序。以以太坊为代表的智能合约平台上，交易排序质量对区块手续费收益、交易执行成功率和用户体验具有直接影响\cite{daian2019flash}。当前主流排序策略采用基于 Gas Price 的贪心方法，按交易手续费从高到低排列以最大化短期收益。该策略实现简单、计算代价低，但存在明显局限：一方面，贪心排列忽略交易之间的状态读写依赖与资源冲突，当两笔交易访问相同合约状态时，顺序不当可能导致执行失败或回滚；另一方面，低手续费交易在该策略下长期得不到确认，等待时间分布容易失衡。

最大可提取价值（Maximal Extractable Value，MEV）问题进一步加剧了上述矛盾\cite{daian2019flash}。MEV 是指区块构建者通过插入、删除或重排交易所能获得的额外利润，其规模在以太坊生态中已达到相当水平。典型的 MEV 攻击形式包括抢跑交易（front-running）、尾随交易（back-running）和三明治攻击（sandwich attack）\cite{chi2024remeasure}。在三明治攻击中，攻击者在目标交易前后分别插入买入和卖出交易，利用目标交易造成的价格变动获利，而目标交易的用户则承受更大的价格滑点。这类行为使普通用户遭受实际经济损失，同时推高了链上 Gas 费用的整体水平。针对 MEV，业界提出了 MEV-Boost\cite{flashbots2022mevboost}、提议者--构建者分离（PBS）\cite{ethereum2024pbs}、公平排序服务和加密内存池等多种机制。然而，这些方案主要在宏观架构层面重塑 MEV 的分配方式或防止交易信息泄露，在构建者内部如何排序交易以平衡收益与风险，仍缺乏系统的算法设计\cite{alipanahloo2024mev}。

近年来，强化学习（Reinforcement Learning，RL）被引入区块链交易排序领域。Wen 等\cite{wen2025rlbes}提出 RL-BES 系统，将深度强化学习与搜索算法结合用于区块构建；也有研究采用深度 Q 网络（DQN）学习排序策略以抑制 MEV。然而，现有 RL 方法在动作空间设计上存在共性不足：其动作空间通常定义为在 FIFO、Gas 优先等少数预定义规则之间切换，本质上是“选规则”而非“逐笔选交易”。这种粗粒度动作设计将排序决策简化为分类任务，难以捕获交易间细粒度依赖关系，也难以根据候选交易集合构成动态调整策略。

基于上述分析，本文提出一种基于深度强化学习的以太坊交易序列决策排序方法。与已有工作的关键区别在于：本文将区块构建过程建模为逐笔交易选择的马尔可夫决策过程（MDP），智能体在每一步从候选交易集合中选取一笔交易加入区块执行序列，直至满足终止条件。这种建模方式将排序决策粒度从“选择排序规则”细化为“选择具体交易”，使策略能够利用当前位置、剩余容量和已选序列上下文做条件决策。本文主要贡献包括三个方面：（1）构建中高负载区块构建场景下的可配置交易排序仿真框架，统一候选池规模、风险比例、Gas 上限与统计评估设置；（2）提出收益--等待时间均衡--风险暴露三目标统一优化的 RL 排序方法，在奖励中引入增量公平项与位置敏感风险惩罚；（3）在统一实验矩阵下，与 FIFO、Gas 优先、启发式风险感知、Fee-Risk 线性评分和 Fair-Fee 双目标贪心五类基线进行系统对比，并通过鲁棒性与消融分析验证方法稳定性与可解释性。


% 第2章 相关工作
% chapters/journal/02-related-work.tex
% 第2章 相关工作

\section{相关工作}

\subsection{传统交易排序策略}

以太坊等区块链平台的默认交易排序遵循 Gas Price 贪心规则：区块构建者按交易手续费从高到低排列候选交易，在区块 Gas 上限约束下依次打包\cite{ethereum2024pbs}。该策略的优化目标单一，仅关注短期手续费收益，对于交易间是否存在状态冲突、执行顺序是否合理等问题不做判断。以太坊 EIP-1559 引入基础费用和优先费用的双层定价机制后，交易手续费的结构更加复杂，简单按优先费用排序并不总能反映交易的真实价值和紧迫程度。

先进先出（FIFO）排序按交易到达时间顺序打包，在公平性方面优于贪心排序，但完全放弃了对收益的优化。\"Oz 等\cite{oz2024fcfs}在先到先得排序机制下研究了 MEV 提取行为，指出即使采用时序排列，攻击者仍可通过网络层面的时延操纵影响交易顺序，FIFO 并不能从根本上消除排序操纵的可能性。Hay 等\cite{hay2025fairness}提出基于乐观策略的交易排序公平性协议，试图在去中心化环境下保证交易排序的时序公平，但其依赖可信时间源假设，在实际部署中面临时钟同步和网络延迟等工程挑战。上述方法均为基于静态规则的排序策略，在设计时预设了固定的排序准则，缺乏对交易间依赖关系和排序风险的动态感知能力。当候选交易集合的构成发生变化时，静态规则无法做出自适应调整。

\subsection{MEV 问题与应对方案}

Daian 等\cite{daian2019flash}在 Flash Boys 2.0 中系统揭示了以太坊 DEX 场景下的 MEV 问题，指出套利机器人通过抢跑和三明治攻击操纵交易排序，从普通用户交易中获取额外收益。Mazorra 等\cite{mazorra2022price}从博弈论视角分析了 MEV 的价格效应。在工程实践方面，Flashbots 提出 MEV-Boost\cite{flashbots2022mevboost}，通过中继将验证者与构建者解耦，但该机制不改变构建者内部的排序逻辑。Vedula 等\cite{vedula2023masquerade}提出 Masquerade 方案，通过轻量级交易重排缓解 MEV，但其重排规则为静态设计，适应性有限。Alipanahloo 等\cite{alipanahloo2024mev}对现有 MEV 缓解方案进行了全面综述，指出在构建者层面缺乏同时兼顾收益与风险控制的排序方法。

\subsection{强化学习在交易排序中的应用}

强化学习在组合优化和序列决策问题上已展现出良好性能\cite{sutton2018rl}。区块构建本质上是在资源约束下从候选集合中构造有序序列的问题：给定一组候选交易和区块 Gas 容量上限，目标是选择交易子集并确定执行顺序，使得特定目标函数最大化。该问题与经典的带约束排列优化具有相似的数学结构，而强化学习的序列决策框架天然适合对这类问题进行建模。

在区块链交易排序领域，已有若干将深度强化学习引入的尝试。Wen 等\cite{wen2025rlbes}提出 RL-BES，采用双 DRL 网络结合蒙特卡洛树搜索进行区块构建，在提升 MEV 提取效率方面取得了效果，但其优化目标侧重于构建者收益最大化，未考虑用户公平性和风险抑制。Seoev 等\cite{seoev2025bidding}将强化学习用于 Polygon 链上的 MEV 提取竞价策略，关注点在于竞价博弈而非排序本身。在区块链其他领域，深度强化学习也被用于共识参数优化\cite{zou2023consensus}和支付通道调度\cite{ren2024payment}等任务，证明了 RL 在区块链系统优化中的普适性。

表~\ref{tab:comparison} 从动作空间设计、优化目标和风险考量三个维度对比了本文方法与已有 RL 排序工作的区别。

\begin{table}[htbp]
\centering
\caption{本文方法与已有 RL 排序方法的对比}
\label{tab:comparison}
\begin{tabular}{lccc}
\toprule
方法 & 动作空间 & 优化目标 & 风险考量 \\
\midrule
RL-BES\cite{wen2025rlbes}  & 规则切换 & MEV 提取效率 & 无 \\
DQN 排序     & 规则切换 & 综合指标 & 有限 \\
本文方法     & 逐笔选交易 & 收益+公平+风险 & 复合奖励 \\
\bottomrule
\end{tabular}
\end{table}

与上述工作相比，本文方法在动作空间设计上具有本质区别。现有 RL 排序方法的动作空间通常定义为在 FIFO、Gas 优先等少量预设规则之间切换，智能体每步选择一种排序规则应用于整个交易集合。这种"选规则"的动作设计将排序决策简化为分类问题，策略网络只能在有限的规则集合中做出选择，无法刻画交易间的位置依赖关系。本文将动作空间定义为"逐笔选交易"：智能体在每步从当前候选集合中选择一笔具体交易，经过多步决策生成完整的排序序列。该设计使策略网络能够感知每笔交易在序列中的位置效应，根据已选交易的构成动态调整后续选择，从而在更细的粒度上优化排序结果。


% 第3章 问题建模
% chapters/journal/03-problem-modeling.tex
% 第3章 问题建模

\section{问题建模}

\subsection{区块构建的序列决策过程}

将一次区块构建过程定义为一个马尔可夫决策过程的完整回合（episode）。给定交易池中的待处理交易集合 $\mathcal{T} = \{tx_1, tx_2, \ldots, tx_N\}$ 和区块 Gas 上限 $G_{\max}$，智能体从空区块开始，在每个时间步 $t$ 从剩余候选交易中选取一笔加入区块执行序列，直至满足终止条件。终止条件包括：候选交易集合为空、区块剩余 Gas 容量不足以容纳任何候选交易、已选交易数量达到预设上限、或智能体主动选择终止动作。回合结束时输出区块内交易的完整执行序列 $\sigma = (tx_{\pi(1)}, tx_{\pi(2)}, \ldots, tx_{\pi(K)})$，其中 $K$ 为实际打包的交易数量，$\pi$ 为排列函数。

该建模方式将区块构建视为一个有限时域的序列决策问题。与将整个候选交易集合一次性排序的方法不同，逐笔选择的建模方式允许智能体在每步决策时利用已选交易序列的信息，使后续选择能够考虑前序交易对区块状态的影响。例如，当一笔高 Gas 消耗的交易被选入后，剩余容量的变化会影响后续可选交易的范围；当一笔与特定合约交互的交易被选入后，其他操作同一合约的交易的执行结果可能因状态变化而不同。这种动态信息在一次性排序方法中难以被充分利用。

\subsection{状态空间}

时间步 $t$ 的状态 $s_t$ 由候选交易特征和区块构建状态两部分组成：

\begin{equation}
  s_t = (\mathbf{F}_t, \mathbf{b}_t)
\end{equation}

候选交易特征矩阵 $\mathbf{F}_t \in \mathbb{R}^{|\mathcal{T}_t| \times d}$ 描述当前剩余候选交易集合 $\mathcal{T}_t$ 中每笔交易的 $d$ 维特征向量。每笔交易 $tx_i$ 的特征向量 $\mathbf{f}_i$ 包含以下分量：交易手续费参数（maxPriorityFeePerGas 或 gasPrice）、Gas 消耗估计 $g_i$、到达时间戳 $t_i^{\text{arr}}$、交易类型编码、以及基于历史数据统计的风险评分 $r_i^{\text{risk}}$。

区块构建状态向量 $\mathbf{b}_t$ 包含：区块剩余 Gas 容量 $g_t^{\text{rem}} = G_{\max} - \sum_{j<t} g_{\pi(j)}$、已选交易数量 $k_t$、已累计手续费收益 $R_t^{\text{fee}}$、以及已选交易序列的特征摘要向量 $\mathbf{h}_t^{\text{seq}}$。其中 $\mathbf{h}_t^{\text{seq}}$ 通过对已选交易特征取均值池化获得。

\subsection{动作空间}

在时间步 $t$，动作空间定义为当前候选交易集合与终止动作的并集：

\begin{equation}
  \mathcal{A}_t = \mathcal{T}_t \cup \{\texttt{STOP}\}
\end{equation}

为保证动作的合法性，对候选交易施加以下过滤规则，生成合法动作集合 $\mathcal{A}_t^{\text{valid}}$：

（1）容量约束：候选交易的 Gas 消耗估计不超过区块剩余容量，即 $g_i \leq g_t^{\text{rem}}$。

（2）依赖约束：候选交易的前置依赖交易（如同一发送地址的 Nonce 顺序）已全部被选入执行序列。

不满足上述约束的交易从候选集合中移除，策略网络仅在合法动作集合上进行决策。

\subsection{奖励函数}

奖励函数设计为手续费收益、确认公平性和风险抑制三项的加权组合。当智能体在时间步 $t$ 选择交易 $tx_{\pi(t)}$ 时，即时奖励为：

\begin{equation}
  r_t = \alpha \cdot r_t^{\text{fee}} + \beta \cdot r_t^{\text{fair}} - \gamma \cdot r_t^{\text{risk}}
\end{equation}

其中 $\alpha, \beta, \gamma$ 为预设权重系数。各分量定义如下：

手续费收益项 $r_t^{\text{fee}}$ 为选中交易的归一化手续费：

\begin{equation}
  r_t^{\text{fee}} = \frac{\text{fee}(tx_{\pi(t)})}{\max_{tx \in \mathcal{T}} \text{fee}(tx)}
\end{equation}

确认公平性项 $r_t^{\text{fair}}$ 衡量交易等待时间的合理性，等待时间越长的交易被选中时给予越高的公平性奖励：

\begin{equation}
  r_t^{\text{fair}} = \frac{t_{\text{now}} - t_{\pi(t)}^{\text{arr}}}{T_{\max}}
\end{equation}

其中 $T_{\max}$ 为归一化常数。该项鼓励策略适当关注长时间未被确认的交易，避免低手续费交易被无限期搁置。

风险惩罚项 $r_t^{\text{risk}}$ 基于交易的风险评分与位置敏感度的乘积：

\begin{equation}
  r_t^{\text{risk}} = r_{\pi(t)}^{\text{risk}} \cdot \phi(t, K)
\end{equation}

其中 $r_{\pi(t)}^{\text{risk}}$ 为交易的风险评分，$\phi(t, K)$ 为位置敏感度函数，在区块头部和尾部位置取较高值（这些位置对抢跑和尾随攻击更为敏感），在中间位置取较低值。具体地，$\phi(t, K)$ 定义为：

\begin{equation}
  \phi(t, K) = 1 - \frac{4t(K-t)}{K^2}
\end{equation}

该函数在 $t=0$（区块首位）和 $t=K$（区块末位）处取值为 1，在 $t=K/2$（区块中间）处取最小值 0，形成 U 形分布。这一设计的动机在于：抢跑攻击通常依赖攻击交易被放置在目标交易之前（即区块靠前位置），而尾随攻击则依赖攻击交易紧随目标交易之后（即区块靠后位置）。将高风险交易排列在区块中间位置可以降低其被利用为攻击工具的可能性。

奖励函数中三个权重系数 $\alpha$、$\beta$、$\gamma$ 控制了收益、公平性和风险抑制之间的相对重要性。在本文的实验配置中，默认取 $\alpha = 1.0$、$\beta = 0.3$、$\gamma = 0.5$。收益项权重最大，保证策略不会过度牺牲区块手续费收入；公平性项和风险项作为辅助目标，引导策略在维持收益水平的同时兼顾确认公平性和排序安全性。


% 第4章 方法设计
% chapters/journal/04-method.tex
% 第4章 方法设计

\section{方法设计}

\subsection{交易特征提取}

对每笔候选交易 $tx_i$，提取 $d$ 维原始特征向量 $\mathbf{x}_i$ 后，通过两层全连接网络编码为稠密表示：

\begin{equation}
  \mathbf{h}_i = \text{ReLU}(\mathbf{W}_2 \cdot \text{ReLU}(\mathbf{W}_1 \mathbf{x}_i + \mathbf{b}_1) + \mathbf{b}_2)
\end{equation}

其中 $\mathbf{h}_i \in \mathbb{R}^{d_h}$ 为交易 $tx_i$ 的编码表示，$d_h$ 为隐藏层维度。交易特征编码网络的参数在策略网络训练过程中联合优化。

原始特征向量 $\mathbf{x}_i$ 由三组特征拼接而成。交易基础特征包括归一化手续费、Gas 消耗估计、到达时间戳和交易类型的独热编码。历史统计特征包括发送地址的历史交易成功率和历史确认延迟。风险特征包括基于历史异常行为统计生成的风险评分，该评分综合考虑发送地址参与异常交易（如失败率过高、高频提交）的历史频率。

\subsection{策略网络结构}

本文采用近端策略优化（Proximal Policy Optimization，PPO）\cite{schulman2017ppo}算法训练排序策略。选择 PPO 而非 DQN\cite{mnih2015dqn} 的原因在于：本文的动作空间随候选交易集合的变化而动态改变，传统 DQN 的固定输出维度难以适应这一特性；PPO 基于策略梯度的框架通过对每笔候选交易独立计算选择得分，天然支持变长动作空间。策略网络采用 Actor-Critic 架构，Actor 和 Critic 共享交易特征编码层，在此基础上分别接各自的决策头。

Actor 网络接收状态 $s_t$，对每笔候选交易计算选择得分：

\begin{equation}
  u_i = \mathbf{w}_a^\top \text{ReLU}(\mathbf{W}_s [\mathbf{h}_i; \mathbf{b}_t] + \mathbf{b}_s)
\end{equation}

其中 $[\mathbf{h}_i; \mathbf{b}_t]$ 为交易表示与区块状态的拼接向量。选择得分经 softmax 归一化后得到动作概率分布 $\pi_\theta(a|s_t)$。

Critic 网络估计当前状态的价值函数 $V_\psi(s_t)$，用于计算优势函数以降低策略梯度的方差。Critic 网络对所有候选交易表示取均值池化后与区块状态拼接，经两层全连接网络输出标量价值估计：

\begin{equation}
  V_\psi(s_t) = \mathbf{w}_v^\top \text{ReLU}(\mathbf{W}_v [\bar{\mathbf{h}}_t; \mathbf{b}_t] + \mathbf{b}_v) + b_v
\end{equation}

其中 $\bar{\mathbf{h}}_t = \frac{1}{|\mathcal{T}_t|}\sum_{tx_i \in \mathcal{T}_t} \mathbf{h}_i$ 为候选交易表示的均值池化。与 Actor 网络需要对每笔交易分别计算得分不同，Critic 网络通过均值池化将变长的候选集合压缩为固定维度的向量，从而输出对当前状态的整体价值估计。

\subsection{合法性掩码机制}

由于候选交易集合的大小在不同时间步动态变化，且部分交易受容量约束和依赖约束限制不可被选取，本文引入合法性掩码（action masking）机制。在计算动作概率分布时，将不属于合法动作集合 $\mathcal{A}_t^{\text{valid}}$ 的交易得分设为负无穷，再进行 softmax 归一化：

\begin{equation}
  \pi_\theta^{\text{mask}}(a_i|s_t) = \frac{\exp(u_i) \cdot \mathbb{I}(a_i \in \mathcal{A}_t^{\text{valid}})}{\sum_{a_j \in \mathcal{A}_t^{\text{valid}}} \exp(u_j)}
\end{equation}

该机制确保策略网络仅在合法动作上分配概率，避免生成违反约束的无效排序。

为支持结构消融实验，本文在实现中进一步提供两类可控开关：其一，关闭动作掩码（No-ActionMask）以检验合法性先验对策略学习的贡献；其二，关闭 \texttt{STOP} 动作（No-STOP）以检验主动终止能力对风险控制和打包效率的影响。完整模型保留动作掩码与 \texttt{STOP}，作为默认配置。

\subsection{训练算法}

训练过程采用 PPO 的截断替代目标函数。在每个训练回合中，智能体与仿真环境交互完成一次完整的区块构建过程，收集轨迹 $\tau = \{(s_t, a_t, r_t, s_{t+1})\}_{t=0}^{T}$。即时奖励按第3章定义计算，即手续费收益、增量等待时间均衡奖励、位置敏感风险惩罚与提前终止惩罚的加权组合。使用广义优势估计（Generalized Advantage Estimation，GAE）计算优势函数：

\begin{equation}
  \hat{A}_t = \sum_{l=0}^{T-t} (\gamma \lambda_{\text{GAE}})^l \delta_{t+l}
\end{equation}

其中 $\delta_t = r_t + \gamma V_\psi(s_{t+1}) - V_\psi(s_t)$ 为时序差分残差，$\gamma$ 为折扣因子，$\lambda_{\text{GAE}}$ 为 GAE 参数。

策略网络的优化目标为 PPO 截断目标：

\begin{equation}
  L^{\text{CLIP}}(\theta) = \mathbb{E}_t \left[ \min\left( \rho_t \hat{A}_t, \; \text{clip}(\rho_t, 1-\epsilon, 1+\epsilon) \hat{A}_t \right) \right]
\end{equation}

其中 $\rho_t = \pi_\theta(a_t|s_t) / \pi_{\theta_{\text{old}}}(a_t|s_t)$ 为重要性采样比率，$\epsilon$ 为截断参数。Critic 网络通过最小化价值函数的均方误差进行更新：

\begin{equation}
  L^{\text{value}}(\psi) = \mathbb{E}_t \left[ (V_\psi(s_t) - \hat{R}_t)^2 \right]
\end{equation}

其中 $\hat{R}_t = \sum_{l=0}^{T-t} \gamma^l r_{t+l}$ 为从时间步 $t$ 开始的折扣累积奖励。训练过程中的主要超参数配置如表~\ref{tab:hyperparams} 所示。

\begin{table}[htbp]
\centering
\caption{训练超参数配置}
\label{tab:hyperparams}
\begin{tabular}{llc}
\toprule
超参数 & 符号 & 取值 \\
\midrule
Actor 学习率 & $\alpha_\theta$ & $3 \times 10^{-4}$ \\
Critic 学习率 & $\alpha_\psi$ & $1 \times 10^{-3}$ \\
折扣因子 & $\gamma$ & 0.99 \\
GAE 参数 & $\lambda_{\text{GAE}}$ & 0.95 \\
PPO 截断参数 & $\epsilon$ & 0.2 \\
每回合更新轮数 & $E$ & 4 \\
隐藏层维度 & $d_h$ & 128 \\
训练回合数 & $M$ & 3{,}000 \\
\bottomrule
\end{tabular}
\end{table}

训练算法的完整流程如算法~\ref{alg:training} 所示。

\begin{algorithm}[H]
\caption{基于 PPO 的交易排序策略训练}
\label{alg:training}
\begin{algorithmic}[1]
\REQUIRE 仿真环境 $\text{Env}$, 学习率 $\alpha_\theta, \alpha_\psi$, 截断参数 $\epsilon$
\ENSURE 训练后的策略参数 $\theta^*$
\STATE 初始化 Actor 参数 $\theta$, Critic 参数 $\psi$
\FOR{episode $= 1, 2, \ldots, M$}
    \STATE 从环境采样候选交易集合 $\mathcal{T}$，初始化区块状态
    \STATE 收集完整轨迹 $\tau = \{(s_t, a_t, r_t)\}_{t=0}^{T}$
    \STATE 计算 GAE 优势估计 $\hat{A}_t$
    \FOR{epoch $= 1, 2, \ldots, E$}
        \STATE 更新 $\theta \leftarrow \theta + \alpha_\theta \nabla_\theta L^{\text{CLIP}}(\theta)$
        \STATE 更新 $\psi \leftarrow \psi - \alpha_\psi \nabla_\psi L^{\text{value}}(\psi)$
    \ENDFOR
\ENDFOR
\RETURN $\theta^*$
\end{algorithmic}
\end{algorithm}

推理阶段，策略网络对候选交易计算选择得分并施加合法性掩码后，取概率最高的交易作为下一步选择（贪心解码），依次构建完整的区块执行序列。

\subsection{推理过程}

训练完成后，策略网络以贪心方式执行推理。给定一组候选交易集合 $\mathcal{T}$ 和初始空区块状态，推理过程按以下步骤逐笔构建区块执行序列：

（1）计算当前时间步所有候选交易的选择得分，并施加合法性掩码过滤不可选交易。

（2）从合法动作集合中选取得分最高的交易 $a_t^* = \arg\max_{a \in \mathcal{A}_t^{\text{valid}}} \pi_\theta^{\text{mask}}(a|s_t)$ 加入区块执行序列。若 \texttt{STOP} 动作得分最高，则终止当前区块构建。

（3）更新区块状态，重建合法动作集合，进入下一时间步。

（4）重复上述过程直至满足终止条件：合法候选交易为空、剩余 Gas 不足、达到最大选择步数、或智能体选择 \texttt{STOP}。

推理结束后输出完整的交易执行序列 $\sigma = (tx_{\pi(1)}, \ldots, tx_{\pi(K)})$。与训练阶段的随机采样不同，推理阶段采用确定性贪心选择，以保证评估结果的可复现性。

\subsection{消融设计对应的结构说明}

为验证模型关键组件的必要性，本文设计四个结构变体用于消融实验：

（1）\textbf{No-SeqSummary}：移除区块状态向量中的已选交易序列摘要 $\mathbf{h}_t^{\text{seq}}$，Actor 和 Critic 网络仅接收交易特征与标量区块状态（剩余 Gas、已选数量、累计收益）作为输入。该变体用于检验序列上下文信息对逐笔决策质量的贡献。

（2）\textbf{No-ActionMask}：关闭合法性掩码机制，策略网络在全部候选交易（含不合法交易）上计算概率分布。当智能体选择不合法交易时，环境不执行该动作并施加负奖励惩罚，但不会因此进入无限循环。该变体用于评估合法动作先验对策略学习效率和最终性能的影响。

（3）\textbf{No-STOP}：移除 \texttt{STOP} 动作，智能体只能通过外部终止条件（候选池耗尽、Gas 不足、达到步数上限）结束区块构建。该变体用于检验主动终止能力对风险控制和打包效率的贡献。

（4）\textbf{Full Model}：保留所有组件（序列摘要、动作掩码、\texttt{STOP} 动作）的完整模型，作为消融实验的对照基准。


% 第5章 实验设计与分析
% chapters/journal/05-experiments.tex
% 第5章 实验设计与分析

\section{实验设计与分析}

\subsection{仿真环境}

本文构建了一个可配置的通用交易排序仿真环境，用于模拟以太坊区块构建过程。与直接使用以太坊测试网或 Geth 私链相比，仿真环境具有两方面优势：一是训练效率高，单次区块构建回合的仿真时间在毫秒级别，可支持大规模并行采样；二是参数可控，能够精确调节候选交易池规模、风险交易比例等关键变量，便于设计对照实验。环境的核心参数如表~\ref{tab:env_params} 所示。

\begin{table}[htbp]
\centering
\caption{仿真环境参数配置}
\label{tab:env_params}
\begin{tabular}{lll}
\toprule
参数 & 符号 & 取值 \\
\midrule
区块 Gas 上限 & $G_{\max}$ & 30{,}000{,}000 \\
候选交易池大小 & $N$ & 50--200 \\
交易 Gas 消耗范围 & $g_i$ & 21{,}000--500{,}000 \\
手续费分布 & -- & 对数正态分布 \\
风险交易比例 & $p_{\text{risk}}$ & 5\%--30\% \\
交易到达时间 & $t_i^{\text{arr}}$ & 泊松过程 \\
\bottomrule
\end{tabular}
\end{table}

仿真环境在每个回合开始时，从预设分布中采样生成候选交易集合。每笔交易的手续费参数服从对数正态分布，以模拟实际以太坊交易手续费的右偏特征。Gas 消耗根据交易类型确定：简单转账交易固定为 21{,}000 Gas，合约调用交易从 $[50{,}000, 500{,}000]$ 区间均匀采样。交易到达时间由泊松过程生成，平均到达间隔设为 1 秒。

风险交易按预设比例 $p_{\text{risk}}$ 注入候选集合。风险交易的风险评分从均匀分布 $U(0.5, 1.0)$ 中采样，普通交易的风险评分从 $U(0, 0.3)$ 中采样，两类交易在风险评分上存在明显的分布差异。风险交易的手续费通常偏高，以模拟实际场景中 MEV 攻击交易通过提高手续费争取优先打包的行为。环境支持调节 $p_{\text{risk}}$ 以模拟不同风险负载场景，从低风险（$p_{\text{risk}} = 5\%$）到高风险（$p_{\text{risk}} = 30\%$）。

交易间的依赖关系通过 Nonce 约束建模：同一发送地址的多笔交易必须按 Nonce 从小到大的顺序被选入执行序列。仿真环境在每个时间步自动检查依赖约束，仅将满足约束的交易保留在合法动作集合中。

\subsection{基线方法与评价指标}

本文选取以下三种基线方法进行对比：

（1）FIFO 排序：按交易到达时间顺序依次打包，直至区块 Gas 用尽。该方法不做任何优化，代表了最朴素的公平排序方案。

（2）Gas 优先排序：按交易手续费从高到低排列并依次打包，为当前以太坊客户端的默认排序策略。该方法追求收益最大化，但不考虑公平性和风险。

（3）启发式风险感知排序（Heuristic）：在 Gas 优先排序基础上，对风险评分超过阈值 $\tau_{\text{risk}} = 0.5$ 的交易降低优先级，将其排列在区块中间位置（即非头尾区域）。该方法代表了基于规则的风险感知排序方案。

评价指标包括以下四项。区块手续费收益（Block Fee Revenue）定义为区块内所有交易手续费之和，以 Gwei 为单位，衡量排序策略对区块构建者的经济效益。确认公平性指数（Fairness Index）采用 Jain 公平性指数衡量交易等待时间分布的均匀程度：

\begin{equation}
  J = \frac{(\sum_{i=1}^K w_i)^2}{K \cdot \sum_{i=1}^K w_i^2}
\end{equation}

其中 $w_i$ 为第 $i$ 笔被打包交易的等待时间，$K$ 为打包交易总数。$J$ 的取值范围为 $[1/K, 1]$，越接近 1 表示等待时间分布越均匀，排序越公平。风险暴露度（Risk Exposure）定义为风险评分超过 $\tau_{\text{risk}}$ 的交易出现在区块前 10\% 和后 10\% 位置的比例，越低表示高风险交易被放置在敏感位置的可能性越小。Gas 利用率（Gas Utilization）定义为区块实际消耗 Gas 总量占区块 Gas 上限 $G_{\max}$ 的比例，反映区块空间的利用效率。

\subsection{主要实验结果}

表~\ref{tab:main_results} 展示了各方法在默认参数配置（$N=100$，$p_{\text{risk}}=15\%$）下的实验结果。

\begin{table}[htbp]
\centering
\caption{各方法在默认配置下的性能对比}
\label{tab:main_results}
\begin{tabular}{lcccc}
\toprule
方法 & 区块收益 & 公平性指数 & 风险暴露度 & Gas 利用率 \\
\midrule
FIFO        & \TODO{} & \TODO{} & \TODO{} & \TODO{} \\
Gas 优先    & \TODO{} & \TODO{} & \TODO{} & \TODO{} \\
Heuristic   & \TODO{} & \TODO{} & \TODO{} & \TODO{} \\
本文方法    & \TODO{} & \TODO{} & \TODO{} & \TODO{} \\
\bottomrule
\end{tabular}
\end{table}

\TODO{在实验数据就绪后填入具体数值，并撰写结果分析段落。分析要点：（1）与 Gas 优先相比，本文方法在收益方面的表现；（2）与 FIFO 相比，公平性指标的差异；（3）与 Heuristic 相比，风险暴露度的改善；（4）Gas 利用率是否保持在合理水平。}

\subsection{鲁棒性分析}

为验证本文方法在不同条件下的稳定性，设计了风险负载变化实验。将风险交易比例 $p_{\text{risk}}$ 分别设为 5\%、10\%、15\%、20\% 和 30\%，在每种配置下运行各方法并记录指标。

\begin{table}[htbp]
\centering
\caption{不同风险负载下各方法的区块收益与风险暴露度}
\label{tab:robustness}
\begin{tabular}{lcccccc}
\toprule
\multirow{2}{*}{方法} & \multicolumn{3}{c}{区块收益} & \multicolumn{3}{c}{风险暴露度} \\
\cmidrule(lr){2-4} \cmidrule(lr){5-7}
 & 5\% & 15\% & 30\% & 5\% & 15\% & 30\% \\
\midrule
FIFO        & \TODO{} & \TODO{} & \TODO{} & \TODO{} & \TODO{} & \TODO{} \\
Gas 优先    & \TODO{} & \TODO{} & \TODO{} & \TODO{} & \TODO{} & \TODO{} \\
Heuristic   & \TODO{} & \TODO{} & \TODO{} & \TODO{} & \TODO{} & \TODO{} \\
本文方法    & \TODO{} & \TODO{} & \TODO{} & \TODO{} & \TODO{} & \TODO{} \\
\bottomrule
\end{tabular}
\end{table}

\TODO{在实验数据就绪后填入具体数值，并撰写鲁棒性分析段落。重点关注：（1）随风险负载增大，各方法性能变化趋势；（2）本文方法在高风险负载下是否仍保持优势；（3）训练收敛曲线的稳定性。}


% 第6章 结论
% chapters/journal/06-conclusion.tex
% 第6章 结论

\section{结论}

本文针对区块链交易排序中现有方法缺乏细粒度序列决策能力的问题，提出了一种基于深度强化学习的以太坊交易序列决策排序方法。该方法将区块构建过程建模为逐笔交易选择的马尔可夫决策过程，在每个时间步从候选交易集合中选取一笔交易加入区块执行序列，通过近端策略优化算法结合合法性掩码机制训练排序策略。奖励函数综合考虑手续费收益、确认公平性和排序风险抑制三个维度，使策略在多个目标之间取得平衡。实验基于可配置的通用交易排序仿真环境，与 FIFO、Gas 优先和启发式风险感知排序三种基线进行对比，验证了本文方法在区块收益、公平性和风险抑制方面的有效性及鲁棒性。

后续工作将在两个方向上扩展：一是将排序环境从通用交易池扩展到去中心化交易所（DEX）场景，引入自动做市商（AMM）价格模型、滑点估计和三明治攻击注入机制；二是将当前的加权奖励设计升级为约束强化学习框架，通过动态拉格朗日乘子实现收益与风险防护的自适应平衡，以应对更复杂的 DeFi 交易环境。


% 参考文献
\bibliographystyle{gbt7714-numerical}
\bibliography{references}

\end{document}
